% Abstract body only. The surrounding \begin{abstract}...\end{abstract}
% (and keyword handling) lives in each build root, since LNCS and the
% two-column preprint wrap the abstract differently.
Real-world data distributions evolve over time, inducing temporal distribution shift that can substantially degrade the reliability of deployed machine learning systems. However, the extent to which architectural choices and their associated inductive biases affect temporal robustness remains insufficiently understood.

We present a systematic empirical comparison of temporal robustness across three heterogeneous, time-indexed domains encompassing image classification, multi-label text classification, and text regression tasks. Using a unified evaluation framework based on temporal drift matrices, we train models on cumulative historical data and evaluate their performance on both earlier and later time periods, thereby quantifying cross-temporal generalization. Our study spans model families ranging from simple multilayer perceptrons and convolutional networks to recurrent networks and pretrained Transformer-based encoders.

Collectively, the results show that architectural inductive biases systematically shape temporal robustness: models whose inductive biases lead them to exploit localized, highly discriminative features attain the highest in-distribution accuracy, yet those features are often the ones that change most over time, so these models degrade fastest, while pretrained encoders that draw on coarser, more stable representations drift more gradually. These observations offer practical guidance for selecting architectures for real-world systems subject to temporal drift.
